\section{System and Experimental Setup}\label{sec:system}

\subsection{Implementation in vLLM}\label{sec:impl}

We implement the system on vLLM 0.27.1 (V1 engine) with three engine extensions and two processes beside the server.

\paragraph{Live feature export.} The trainer needs the target's auxiliary hidden states at every accepted position. Inside the EAGLE-3 path of the model runner, we copy them from the GPU to pinned host memory without blocking the step, and a background thread writes them per request into a shared-memory ring buffer. Rejected positions are dropped, only tensor-parallel rank 0 writes, and a full buffer drops one step instead of stalling serving. Prefix caching has to be disabled, because positions served from the cache have no hidden states.

\paragraph{Co-located trainer.} A separate process on the same GPU trains the drafter with cross-entropy against the target's own next tokens over the drafter's vocabulary (AdamW; \cref{app:manifest}). It reads completed requests from the ring buffer into a last-in-first-out buffer of at most 600k tokens, discards data older than 60\,s, and implements duty cycle $d$ by sleeping $(1-d)/d$ times the duration of each step.

\paragraph{Drafter hot swap and speculation switch.} A collective RPC loads new drafter weights into all workers without restarting the engine; in the frozen runs a swap took a median of \nv{swap.m1.med}\,s for the Qwen3-8B drafter and \nv{swap.m2.med}\,s for the MoE drafter. A second RPC switches speculation in the scheduler and the workers together, so drafting stops instead of being computed and discarded; the drafter stays in GPU memory and the round trip takes a median of \nv{switch.med}\,ms.

\paragraph{Controller process.} The controller is a 1\,Hz process that reads the engine's metrics endpoint and the trainer's log, writes the duty cycle to a file that the trainer reads before every step, and requests exports through a command file.

\subsection{Experimental protocol}\label{sec:setup}

\paragraph{One frozen checkpoint.} Every experiment runs on one checkpoint that we froze before the first run: the controller, trainer, and harness code, the six patched engine files, the two calibration files, the factory drafters, the target weights, the prompt pools, the seeds, and the software environment, each recorded by hash (\cref{app:manifest}). \pol{TIDE} was added later as new files only and, like the reruns below, runs in an environment rebuilt from the frozen package list.

\paragraph{Workload.} Each run serves three phases on one H100 GPU: 200 chat requests at concurrency 16 (phase A1), then 16{,}000 unique requests from a new domain at concurrency 256, which saturate the GPU for \nv{bdur.min}--\nv{bdur.max}\,s (phase B), then 200 different chat requests (phase A2). The scenarios are Qwen3-8B with code or Chinese in B and Qwen3-30B-A3B with code in B. Decoding is greedy, outputs stop naturally or at 512 tokens, and the B prompts were never seen during calibration. All policies except \pol{Production} reserve the trainer's memory (\cref{app:manifest}).

\paragraph{Policies.} \pol{Production} runs the factory drafter with speculation always on, without export or memory reserve, as deployed today. \pol{Static} does the same under the evaluation's engine configuration and is the reference for all differences. \pol{Plain} serves without a drafter. \pol{Switch} applies \cref{eq:switch} with the domain's frontier and never trains. \pol{Threshold} turns speculation off whenever the server is saturated and trains at duty 0.5 whenever the predicted TPOT is within the SLO, i.e., the idle-time rule of OSD \citep{liu2024osd} combined with a load threshold. \pol{TIDE} reimplements the training triggers of TIDE \citep{park2026tide} with the speculation rule of \pol{Switch}: it trains at duty 0.5 from the start and whenever acceptance drifts down, swaps in a drafter every 200 requests learned, and restores the previous one and stops if acceptance did not rise. \pol{Always} trains at duty 0.5 from the first request and swaps drafters after 100, 200, 400, \dots{} requests. \pol{Informed} trains at duty 0.5 from the moment phase B starts until it has learned from 800 requests, swapping at 100, 200, 400, and 800; its stopping point comes from the learning curves of \cref{fig:motivation}c, so it is a reference that knows when to stop, not a deployable policy. \pol{Full} is the controller of \cref{sec:method}. Neither turns speculation off (\cref{app:policies}).

\paragraph{Repetitions and metrics.} Every policy runs three times per scenario. Repetition $r$ fixes the prompt order of all phases and the trainer's random seed, so within a repetition all policies serve the same request sequence, and we report paired differences as mean $\pm$ standard deviation over the three pairs. The main metric is phase-B output throughput; we also report TPOT and whole-run throughput. The \nv{n.runs} runs used \nv{hosts} nodes of the same type and \nv{gpuh} GPU-hours. A scan with a rule fixed in advance found node interference in \nv{scan.flagged.word} of the \nv{scan.n} runs, all of them baselines; we report their reruns, which can only narrow the controller's margins (\cref{app:rerun}).

\paragraph{Further conditions.} Four ablations remove one of \pol{Full}'s tests at a time (payback rule, stop tests, speculation switch, or label and learning curve) on the same checkpoint, prompts, and seeds (\cref{app:ablation}). Two mixed-traffic scenarios serve a 50/50 mixture of code or Chinese with chat in phase B, calibrated as domains of their own, with \pol{Full} also run under a stale label and without any label (\cref{app:mix}). Five policies were rerun on a port to SGLang 0.5.10 with its own frozen calibration (\cref{app:sglang}), and a fluid simulator built from the frozen calibration replays unseen traces against four criteria declared before the runs (\cref{app:sim}).
